% ---------------------------------------------------------------------------
% Shared preamble: palette, chess-diagram macros, callout boxes.
% Compile the main file with lualatex (needed for the Unicode chess glyphs).
% ---------------------------------------------------------------------------
% Fail here, naming the engine and the command, rather than fifty lines later
% inside fontspec.  This is the first thing the preamble does.
\RequirePackage{iftex}
\newcommand{\kqkNeedsLuaTeX}{%
  \PackageError{endgames}{lualatex is required to build this document}%
    {The chess men are Unicode glyphs from a system font, and every diagram is
     built\MessageBreak by an embedded Lua module, so pdflatex and xelatex
     cannot compile it.\MessageBreak Run `make' in doc/, or `lualatex
     endgames.tex' twice for the contents.}%
  \stop}
% The name may not begin with \end: LaTeX reserves that prefix for environment
% terminators and \newcommand refuses it.  \expandafter closes the conditional
% before \stop fires, so the run ends on the message above rather than on an
% incomplete \iffalse.
\ifLuaTeX\else\expandafter\kqkNeedsLuaTeX\fi
\usepackage[a4paper,margin=22mm,top=20mm,bottom=20mm]{geometry}
\usepackage{fontspec}
\usepackage{microtype}
\usepackage{tikz}
\usepackage{pgfplots}
\usepackage{booktabs}
\usepackage{tcolorbox}
\usepackage{amsmath}
\usepackage{amssymb}
\usepackage{titlesec}
\usepackage{needspace}
\usepackage[hidelinks]{hyperref}
% stepboard draws every chess diagram in this document.  Where a diagram has a
% line to follow, the plies are optional-content layers stacked in place and a
% click on the controls swaps which one is shown; where it does not, it is an
% ordinary diagram with the same badges and geometry.  It loads ocgx2 and must
% come after hyperref.  No JavaScript is involved, so a viewer that ignores
% optional content simply shows the opening position.
\usepackage{stepboard}
\pgfplotsset{compat=1.18}
\usetikzlibrary{calc,positioning,decorations.pathreplacing,arrows.meta}
\tcbuselibrary{skins,breakable}

% ---- palette --------------------------------------------------------------
\definecolor{ink}      {HTML}{16202A}
\definecolor{inksoft}  {HTML}{50606E}
\definecolor{paper}    {HTML}{FBF8F3}
\definecolor{kqk}      {HTML}{6C3FA8}   % queen  -- violet
\definecolor{krk}      {HTML}{0E7C7B}   % rook   -- teal
\definecolor{kbbk}     {HTML}{D67200}   % bishops-- amber
\definecolor{kbnk}     {HTML}{C1264B}   % B+N    -- crimson
\definecolor{kqkr}     {HTML}{2F6B3A}   % Q v R  -- forest; the armed defender
\definecolor{kqkk}     {HTML}{17539E}   % Q v KK -- blue; the doubled defender
\definecolor{knnnk}    {HTML}{B5179E}   % 3N v K -- magenta; the knights
\definecolor{knnk}     {HTML}{7A2E93}   % 2N v K -- plum; the table beneath it
\definecolor{kings}    {HTML}{A5452A}   % KKK v KK -- brick; nothing but kings
% Board and men follow the Lichess default: the brown board, white men white
% and black men black.  Nothing about a diagram should have to be decoded --
% colour means side, and only side.
\definecolor{sqlight}  {HTML}{F0D9B5}
\definecolor{sqdark}   {HTML}{B58863}
\definecolor{manwhite} {HTML}{FFFFFF}
\definecolor{manblack} {HTML}{262421}
\definecolor{arrowW}   {HTML}{15803D}   % White's move -- green, as Lichess draws it
\definecolor{arrowB}   {HTML}{B4232C}   % Black's move
\definecolor{hlwin}    {HTML}{86C08A}
\definecolor{hldraw}   {HTML}{E4A0A0}
\definecolor{hlmark}   {HTML}{F2D061}

\newfontfamily\chessfont{Apple Symbols}

% ---- headings -------------------------------------------------------------
\titleformat{\section}{\Large\bfseries\color{ink}}{\thesection}{0.6em}{}[{\color{ink!25}\titlerule[1pt]}]
\titleformat{\subsection}{\large\bfseries\color{inksoft}}{\thesubsection}{0.6em}{}
% Level three sits under level two, so it must not weigh more: same colour,
% one size down.  Without this it inherits the default black and reads as the
% most prominent heading on the page.
\titleformat{\subsubsection}{\normalsize\bfseries\color{inksoft}}{\thesubsubsection}{0.6em}{}
\titlespacing*{\section}      {0pt}{2.6ex plus .6ex minus .2ex}{1.4ex plus .2ex}
\titlespacing*{\subsection}   {0pt}{2.0ex plus .5ex minus .2ex}{0.9ex plus .2ex}
\titlespacing*{\subsubsection}{0pt}{1.5ex plus .4ex minus .2ex}{0.7ex plus .2ex}


% \S\ref{sec:kings} and \S\ref{sec:rules} go one level deeper than the rest, so
% level four is numbered too -- otherwise a \ref to one of those blocks would
% silently resolve to its parent.  Run-in, so it never reads as a heading.
\setcounter{secnumdepth}{4}
\titleformat{\paragraph}[runin]{\normalsize\bfseries\color{inksoft}}{\theparagraph}{0.5em}{}[.]
\titlespacing*{\paragraph}{0pt}{1.4ex plus 1ex minus .2ex}{0.8em}
% ---- chess diagram --------------------------------------------------------
% Squares are addressed by 0-based (file, rank).  \bd{<sq size in cm>}{<n>}{body}
\newcommand{\glyphK}{♔}\newcommand{\glyphQ}{♕}\newcommand{\glyphR}{♖}
\newcommand{\glyphB}{♗}\newcommand{\glyphN}{♘}\newcommand{\glyphP}{♙}
% Lower case is Black, as \glyphk already is.  The stalemate classification in
% \S8 needs a black rook, bishop and pawn; nothing else in the document does.
\newcommand{\glyphk}{♚}\newcommand{\glyphq}{♛}\newcommand{\glyphr}{♜}
\newcommand{\glyphb}{♝}\newcommand{\glyphn}{♞}\newcommand{\glyphp}{♟}

\newlength{\sqsz}\newlength{\glyphsz}

% \bdw{<total width>}{<n>}{body}  -- fixed total width, squares shrink with n.
% This is the only entry point; stepboard.lua calls it for every diagram.
\newcommand{\bdw}[3]{\setlength{\sqsz}{\dimexpr#1/#2\relax}\bdX{#2}{#3}}
\newcommand{\bdX}[2]{%
  \setlength{\glyphsz}{0.78\sqsz}%
  \begin{tikzpicture}[x=\sqsz,y=\sqsz,every node/.style={inner sep=0pt}]
    \foreach \f in {0,...,\the\numexpr#1-1\relax}{%
      \foreach \r in {0,...,\the\numexpr#1-1\relax}{%
        \pgfmathparse{int(mod(\f+\r,2))}%
        \ifnum\pgfmathresult=0 \fill[sqdark] (\f,\r) rectangle ++(1,1);
        \else \fill[sqlight] (\f,\r) rectangle ++(1,1);\fi}}%
    #2
    \draw[ink!70,line width=0.7pt] (0,0) rectangle (#1,#1);
  \end{tikzpicture}}

% Pieces sit on a disc, so they stay visible even when the squares are tiny --
% at n = 60 a square is under a millimetre.  The disc carries the piece's
% colour and the rim contrasts with it, so a man reads as white or black at any
% size and on either shade of square.
\newcommand{\disc}[4]{%
  \fill[#3] (#1+0.5,#2+0.5) circle (0.40);
  \draw[#4,line width=0.4pt] (#1+0.5,#2+0.5) circle (0.40);}
\newcommand{\pw}[3]{\disc{#1}{#2}{manwhite}{manblack!75}%
  \node at (#1+0.5,#2+0.5)
    {\chessfont\fontsize{\glyphsz}{\glyphsz}\selectfont\textcolor{manblack}{#3}};}
% \pbx puts any black man on the board.
\newcommand{\pbx}[3]{\disc{#1}{#2}{manblack}{manwhite!85}%
  \node at (#1+0.5,#2+0.5)
    {\chessfont\fontsize{\glyphsz}{\glyphsz}\selectfont\textcolor{manwhite}{#3}};}

% \coordsn{<n>} -- file letters and rank numbers around a board of any size,
% which the diagrams need since they run from 3x3 up to 22x22.  The file letter
% comes from \alph via a counter, since \foreach hands us the file as a number.
\newcounter{coordfile}
\newcommand{\coordn}[1]{\setcounter{coordfile}{\numexpr#1+1\relax}\alph{coordfile}}
\newcommand{\coordsn}[1]{%
  \foreach \f in {0,...,\the\numexpr#1-1\relax}{%
    \node[below=0.5pt,font=\tiny,color=inksoft] at (\f+0.5,0) {\coordn{\f}};}%
  \foreach \r in {0,...,\the\numexpr#1-1\relax}{%
    \node[left=1pt,font=\tiny,color=inksoft] at (0,\r+0.5)
      {\the\numexpr\r+1\relax};}}

% The decorations below are passed to stepboard's under= and over= keys: under=
% is drawn before the men (tints, rays, the long diagonal), over= after them
% (rings, arrows).
%
% \arw{f0}{r0}{f1}{r1}{colour}{number} -- a numbered move arrow, for over=.
% Drawn over the men rather than under them, because a king move is one square
% long and an arrow shortened clear of both discs would have nothing left of
% it.  The number sits at the midpoint.
\newcommand{\arw}[6]{%
  \draw[#5,opacity=0.82,line width=1.3pt,-{Stealth[length=2.4mm,width=2mm]},
        shorten >=1.2pt,shorten <=1.2pt] (#1+0.5,#2+0.5) -- (#3+0.5,#4+0.5);
  \node[circle,fill=#5,text=white,inner sep=0pt,minimum size=3.1mm,
        font=\tiny\bfseries] at ($(#1+0.5,#2+0.5)!0.5!(#3+0.5,#4+0.5)$) {#6};}

% highlight a square / ring a square / the long diagonal / a confinement box
\newcommand{\hl}[3]{\fill[#3,opacity=0.85] (#1,#2) rectangle ++(1,1);}
\newcommand{\ring}[3]{\draw[#3,line width=1.1pt] (#1+0.06,#2+0.06) rectangle ++(0.88,0.88);}
\newcommand{\diag}[2]{\draw[#2,line width=0.7pt,dash pattern=on 2pt off 2pt,opacity=0.8]
  (0,0) -- (#1,#1);}
\newcommand{\boxoff}[3]{\fill[#3,opacity=0.22] (#1,#1) rectangle (#2,#2);}
% \cagebox{x0}{y0}{x1}{y1}{colour} -- tint a rectangle of squares, 0-based,
% both corners inclusive.  Used to draw the queen's cage.
\newcommand{\cagebox}[5]{\fill[#5,opacity=0.20] (#1,#2) rectangle ({#3+1},{#4+1});}
% \qcut{file}{rank}{n}{colour} -- the queen's file and rank: the two lines the
% black king cannot cross, and hence the walls of the cage.
% \fray{f0}{r0}{f1}{r1}{colour} -- one of the queen's lines of attack, centre to
% centre.  Call it before the piece macros so the discs are drawn over it.
\newcommand{\fray}[5]{%
  \draw[#5,line width=1pt,dash pattern=on 2.5pt off 2pt,opacity=0.95]
    (#1+0.5,#2+0.5) -- (#3+0.5,#4+0.5);}
\newcommand{\qcut}[4]{%
  \draw[#4,line width=0.9pt,dash pattern=on 2.5pt off 2pt,opacity=0.95]
    (#1+0.5,0) -- (#1+0.5,#3) (0,#2+0.5) -- (#3,#2+0.5);}


% ---- callout --------------------------------------------------------------
\newtcolorbox{result}[2][]{enhanced,breakable,colback=#2!5,colframe=#2,
  boxrule=0pt,leftrule=2.5pt,arc=1pt,left=6pt,right=6pt,top=4pt,bottom=4pt,
  before skip=4pt, after skip=4pt,
  fonttitle=\bfseries\color{white},coltitle=white,
  attach boxed title to top left={xshift=5pt,yshift=-2.5pt},
  boxed title style={colback=#2,arc=1pt,boxrule=0pt},#1}
